\section{Impact Statements}\label{sec:impacts}

The shift from centralized to distributed training of large models, may introduce new technical and societal challenges and have the potential to fundamentally reshape the AI landscape. 


\subsection{AI Monopoly and Democratization}\label{subsec:ai_monopoly_and_democratization}


The current AI landscape is characterized by significant concentration of power among a few tech giants, primarily due to their monopoly over massive computing resources and data centers \cite{bommasani2021opportunities}. This monopolistic trend has intensified with companies like OpenAI increasingly moving towards closed systems. 
While open-source alternatives like Llama \cite{touvron2023llama3}, Deepseek \cite{liu2024deepseek} and other community-driven models have made strides towards democratization by releasing model parameters and technical reports \cite{welsh2024democratising}, the gap in computational resources and data access between major AI companies and other players remains substantial and continues to widen. 
This disparity in resources allows tech giants to maintain their absolute dominance in determining the direction of AI development, raising concerns about AI democratization. 


Edge device-based collaborative training presents a promising pathway to democratize AI development \cite{yao2022edge}.
By leveraging the collective computing power of millions of edge devices, this approach could effectively challenge the existing monopolistic structure \cite{dong2025beyond}. 
This democratization of AI training through edge devices could fundamentally reshape the structure of responsibilities and authorities.
If everyone can participate in training LLMs, the AI landscape could fundamentally change. Training decisions would shift from companies to communities, creating shared responsibility for model development \cite{zeng2023distributed}. Global participation would help models reflect diverse cultural perspectives, while allowing communities to adapt models for their local needs.
Furthermore, this decentralized approach could foster a more competitive and innovative AI ecosystem. When the barriers to entry for AI model training are lowered, we can expect to see a broader range of specialized models emerging (like \cite{song2025injecting,MedicalGPT,li2023large,yao2024lawyer,AlphaEvolve}), better suited to local needs and diverse use cases.


\subsection{Fairness and Incentive Mechanisms}\label{subsec:fairness_and_incentive_mechanisms}

The distributed training paradigm introduces new considerations for model fairness and bias mitigation \cite{kheya2024pursuit,li2020fair}. When training occurs across diverse edge devices, the resulting models can potentially better reflect the heterogeneous nature of user populations \cite{wang2020optimizing}. However, this approach also raises concerns about participation bias, where differences in device capabilities or user engagement could lead to underrepresentation of certain groups \cite{kairouz2021advances}. To address these challenges, researchers have proposed various fairness-aware federated learning algorithms \cite{mohri2019agnostic} that aim to ensure equitable model performance across different demographic groups and device types \cite{li2020fair}.

To sustain a distributed training ecosystem, effective incentive mechanisms are crucial for motivating user participation \cite{kang2019incentive}. Traditional approaches like computational resource sharing \cite{khan2020federated} and privacy-preserving reward systems have shown promise in encouraging user engagement. More innovative solutions include token-based reward systems \cite{wang2023incentive} and reputation mechanisms \cite{yang2019federated} that compensate users for their contributions while maintaining system integrity. These incentive structures not only encourage consistent participation but also help ensure the quality of contributed training data \cite{sharghivand2020edge}, creating a sustainable ecosystem for collaborative AI development.

\subsection{Carbon Footprint and Energy Efficiency}\label{subsec:carbon_footprint_and_energy_efficiency}
The shift from centralized to distributed training offers compelling environmental benefits \cite{guidi2024environmental}. Traditional data centers housing large language models face significant energy challenges\cite{sarkar2024carbon} - their high-performance GPUs require extensive cooling systems that consume 30-40\% of total energy \cite{cooling}. In contrast, FL distributes computation across edge devices like smartphones and tablets that operate at much lower temperatures and power levels, eliminating industrial cooling needs \cite{fl_vs_centralized}.

FL also dramatically reduces data transmission energy costs. While centralized approaches require raw data transfer from millions of devices, FL only transmits lightweight model updates, substantially decreasing network energy overhead \cite{fl_data_transmission}. The hardware efficiency gap is striking - edge devices like the NVIDIA Tegra X2 consume just 7.5W during training compared to 250W for data center GPUs \cite{hardware_power}, translating to major carbon footprint reductions, particularly for simpler models \cite{fl_iid}.
By reducing reliance on power-hungry data centers and leveraging existing consumer devices, FL enables more sustainable AI development through optimized energy efficiency and minimized infrastructure needs. This combination of reduced cooling requirements, efficient hardware utilization, and optimized data handling makes FL an environmentally responsible choice for the future of AI training \cite{fl_environmental}. As climate impact becomes increasingly critical, FL's sustainability advantages position it as a key technology for green AI development.
